\documentclass{report}
\usepackage{amsmath,amssymb}
\setlength{\parindent}{0mm}
\setlength{\parskip}{1em}
\begin{document}
\begin{center}
\rule{6in}{1pt} \
{\large Siegfried Cools \\
{\bf An Efficient Multigrid Calculation of the Far Field Map for Helmholtz Equations}}

Department of Mathematics and Computer Science \\ Universiteit Antwerpen \\ Middelheimlaan 1 \\ 2020 Antwerp \\ Belgium
\\
{\tt siegfried.cools@ua.ac.be}\\
Bram Reps\\
Wim Vanroose\end{center}

\textbf{Introduction.}

In this work we present a new efficient calculation method for the far
field amplitude patterns that arise from scattering problems governed by
the $d$-dimensional Helmholtz equation. The calculation of the far field
mapping is typically a two step process. First a Helmholtz problem with
absorbing boundary conditions (PML, ECS) is solved on a finite numerical
box covering the object of interest. In the second step a volume integral
over the Green's function involving the numerical solution yields the
angular dependency of the far field amplitude. The main computational
bottleneck generally lies within the first step, since it requires a
suitable (iterative) method for the solution of a high dimensional
indefinite Helmholtz system.

Preconditioned Krylov subspace methods are currently among the most
efficient numerical algorithms for the solution of high dimensional
positive definite systems. A generalization of this approach led to the
development of the Complex Shifted Laplacian (CSL) preconditioner,
proposed in [1] as an effective Krylov subspace method preconditioner for
Helmholtz problems. The key idea behind CSL is the formulation of a
perturbed Helmholtz problem which includes a complex valued wave number.
This implies a damping in the problem, thus making the preconditioning
system solvable using multigrid in contrast to the original Helmholtz
problem. Recently a variation on CSL by the name of Complex Stretched
Grid (CSG) was proposed [2], introducing a complex valued grid distance
(complex rotation) in the preconditioner.

The far field map computation proposed in this work reformulates the
standard real-valued Green's function integral expression for the far
field amplitude on a complex contour. Hence, one requires the solution of
the Helmholtz equation on a complex contour. It is shown that the latter
problem is equivalent to a CSL problem that can be solved very
efficiently using a multigrid method. However, whereas CSL was previously
only used as a preconditioner, the proposed complex-valued far field map
calculation effectively allows for multigrid to be used as a solver.

\textbf{The far field map for Helmholtz problems.}

The Helmholtz equation is a mathematical representation of the physics
behind a wave scattering at an object $O$ located within a domain $\Omega
\subset \mathbb{R}^d$. The equation is
\begin{equation} \label{eq:helmholtz_in}
\left(-\Delta - k^2(\boldsymbol{x})\right) u(\boldsymbol{x}) = k_0^2
\chi(\boldsymbol{x})u_{in}(\boldsymbol{x}), \quad \boldsymbol{x} \in
\Omega,
\end{equation}
where $\chi(\boldsymbol{x}) = (k^2(\boldsymbol{x})-k_0^2)/k_0^2$
represents the object of interest. Note that $\chi(\boldsymbol{x}) = 0$
for $\boldsymbol{x} \in \Omega \setminus O$. The above equation can in
principle be solved in a numerical box (i.e.~a discretized subset of
$\Omega$) covering the support of $\chi$, with absorbing boundary
conditions along all boundaries, yielding a numerical solution denoted by
$u^N$. The far field scattered wave then satisfies the following
inhomogeneous Helmholtz equation with constant wave number
\begin{equation} \label{eq:HHR}
\left(-\Delta - k_0^2\right) \, u(\boldsymbol{x}) = g(\boldsymbol{x}),
\quad \boldsymbol{x} \in \mathbb{R}^d,
\end{equation}
where $g(\boldsymbol{x})=k_0^2 \chi(\boldsymbol{x})
(u_{in}(\boldsymbol{x})+u^N(\boldsymbol{x}))$. The analytical solution to
the above equation is given by the Green's integral
\begin{equation} \label{eq:scint}
u(\boldsymbol{x}) = \int_{\Omega} G(\boldsymbol{x},
\boldsymbol{x}^\prime) \, g(\boldsymbol{x}^\prime) \, d\boldsymbol{x}',
\quad \boldsymbol{x} \in \mathbb{R}^d.
\end{equation}
Consequently, the asymptotic form of the scattered wave in the direction
of the unit vector $\boldsymbol\alpha \in \mathbb{R}^d$ is
\begin{equation} \label{eq:FFusc}
\lim_{\rho \to \infty} u(\rho,\boldsymbol\alpha) = D(\rho)
F(\boldsymbol\alpha), \qquad \boldsymbol\alpha \in \mathbb{R}^{d},
\end{equation}
where the far field (amplitude) mapping is given by
\begin{equation} \label{eq:FFmap}
F(\boldsymbol \alpha) = \int_\Omega e^{-ik_0 \boldsymbol{x}^\prime \cdot
\boldsymbol \alpha} g(\boldsymbol{x}^\prime) \, d\boldsymbol{x}^\prime.
\end{equation}

\textbf{Calculation on a complex contour.}

The above far field integral expression can be split into a sum of two
contributions $F(\boldsymbol \alpha) = I_1 + I_2$, where
\begin{align}
I_1 &= \int_\Omega e^{-i k_0 \boldsymbol{x} \cdot \boldsymbol \alpha
}\chi(\boldsymbol{x})u_{in}(\boldsymbol{x}) d\boldsymbol{x}, \\
I_2 &= \int_\Omega e^{-i k_0 \boldsymbol{x} \cdot \boldsymbol \alpha }
\chi(\boldsymbol{x}) u^N(\boldsymbol{x}) d\boldsymbol{x}.
\end{align}
Calculation of first integral $I_1$ is generally easy, since it only
requires the expression for the incoming wave. The second integral
however requires the solution $u^N$ of the Helmholtz equation on the
numerical box, which is known to be notoriously hard to obtain using
present-day iterative methods. However, if both $u$ and $\chi$ are
analytical functions, the integral $I_2$ can be calculated over a complex
contour defined by the rotated real domain $Z_1 = \{\boldsymbol{z} \in
\mathbb{C} \,|\, \boldsymbol{z} = \boldsymbol{x} e^{i\gamma} :
\boldsymbol{x} \in \Omega \}$, where $\gamma$ is a fixed rotation angle,
followed by the curved segment $Z_2 = \{\boldsymbol{z} \in \mathbb{C}
\,|\, \boldsymbol{z} = \boldsymbol{b} e^{i\theta} : \boldsymbol{b} \in
\partial\Omega,~ 0 \leq \theta \leq \gamma\}$, closing the contour. The
integral $I_2$ can then be written as
\begin{equation}
I_2 = \int_{Z_1} e^{-i k_0 \boldsymbol{z} \cdot \boldsymbol \alpha
}\chi(\boldsymbol{z})u^N(\boldsymbol{z}) d\boldsymbol{z},
\end{equation}
where the integral over $Z_2$ has vanished because $\chi$ is per
definition zero everywhere outside $O$, thus notably in all points of
$Z_2$. The advantage of this approach is that we only need the value of
$u^N$ evaluated along this complex contour, where the original Helmholtz
problem is reduced to a damped equation. This problem is equivalent to a
CSL problem, and can thus be solved very efficiently using multigrid,
resulting in a fast and scalable computation of the far field mapping.
The new approach is successfully validated on model problems in two and
three spatial dimensions, where its effectiveness and scalability are
demonstrated.

\textbf{References.}

[1] Y.A. Erlangga, C.W. Oosterlee and C. Vuik, \textit{On a class of
preconditioners for solving the {H}elmholtz equation}, App. Num. Math.,
\textbf{50}(3-4), 2004, pp. 409--425.

[2] B. Reps, W. Vanroose and H. bin Zubair, \textit{On the indefinite
{H}elmholtz equation: {C}omplex stretched absorbing boundary layers,
iterative analysis, and preconditioning}, J. Comput. Phys.,
\textbf{229}(22), 2010, pp. 8384--8405.

[3] S. Cools, B. Reps and W. Vanroose, \textit{An efficient multigrid
calculation of the far field map for Helmholtz problems},
arXiv:1211.4461.


\end{document}
